\BourbakiExercisesSection

\begin{enumerate}
\def\labelenumi{\arabic{enumi}.}
\item
  If M is a monogenous A-module, then \(\mathsf{S}(\mathsf{M}) = \mathsf{T}(\mathsf{M})\) . Deduce an example of an injective homomorphism \(j: N \to M\) such that \(\mathsf{S}(j): \mathsf{S}(\mathsf{N}) \to \mathsf{S}(\mathsf{M})\) is not injective ( \(\S 5\) , Exercise 1).
\item
  Establish for symmetric algebras the properties analogous to those of Exercise 8(a) and (b) of § 5.
\item
  Let \(\mathbf{K}\) be a field, \(\mathbf{B}\) the polynomial ring \(\mathbf{K}[\mathbf{X}_1, \mathbf{X}_2]\) and \(\mathfrak{p}\) the ideal of \(\mathbf{B}\) generated by the polynomial \(\mathbf{X}_1^3 - \mathbf{X}_2^2\) .
\end{enumerate}

\begin{enumerate}
\def\labelenumi{(\alph{enumi})}
\item
  Show that p is prime and therefore A = B/p is an integral domain. Let a and b denote the canonical images of \(X_{1}\) and \(X_{2}\) in A; a and b are not invertible in A and \(a^{3} = b^{2}\) .
\item
  Let m be the (maximal) ideal of A generated by a and b. Show that the symmetric algebra \(\mathsf{S}(\mathsf{m})\) is not an integral domain. (Observe that m is identified with the quotient module \(A^{2}/N\) , where N is the submodule of \(A^{2}\) generated by \((b, -a)\) ; conclude that \(\mathsf{S}(\mathsf{m})\) is isomorphic to the quotient ring \(\mathbf{A}[\mathbf{U}, \mathbf{V}]/(b\mathbf{U} - a\mathbf{V})\) and show that in the polynomial ring \(\mathbf{A}[\mathbf{U}, \mathbf{V}]\) the principal ideal \((b\mathbf{V} - a\mathbf{U})\) is not prime, by considering the product \(b^{2}(b\mathbf{V}^{3} - \mathbf{U}^{3})\) .)
\end{enumerate}

¶ 4. Let A be a commutative ring and a an ideal of A. The Rees algebra of the ideal a is the subalgebra R(a) of the polynomial algebra A{[}T{]} in one indeterminate, consisting of the polynomials \(c_{0} + c_{1}\mathrm{T} + \cdots + c_{n}\mathrm{T}^{n}\) , where \(c_{0} \in \mathbf{A}\) and \(c_{k} \in \mathfrak{a}^{k}\) for every integer \(k \geqslant 1\) .

\begin{enumerate}
\def\labelenumi{(\alph{enumi})}
\item
  Show that there exists one and only one surjective A-algebra homomorphism \(r: \mathsf{S}(\mathfrak{a}) \to \mathsf{R}(\mathfrak{a})\) such that the restriction of \(r\) to \(\mathfrak{a}\) is the A-linear mapping \(x \to x\mathbb{T}\) .
\item
  Suppose that the ring A is an integral domain. Let \(a_1, \ldots, a_n\) be elements of A with \(a_n \neq 0\) ; consider the A-algebra homomorphism
\end{enumerate}

\[
u: \mathrm{A} [ \mathrm{X} _ {1}, \dots , \mathrm{X} _ {n} ] \rightarrow \mathrm{A} [ \mathrm{T} ]
\]

such that \(u(\mathbf{X}_i) = a_i\mathrm{T}\) . The kernel \(\mathfrak{r} = \operatorname{Ker}(u)\) is a graded ideal whose homogeneous component \(\mathfrak{r}_m\) of degree \(m\) consists of the homogeneous polynomials \(f\) of degree \(m\) such that \(f(a_1, \ldots, a_n) = 0\) . If \(\mathfrak{r}' \subset \mathfrak{r}\) is the ideal of \(\mathrm{A}[\mathrm{X}_1, \ldots, \mathrm{X}_n]\) generated by the homogeneous component \(\mathfrak{r}_1\) of \(\mathfrak{r}\) , show that the A-module \(\mathfrak{r}/\mathfrak{r}'\) is a torsion A-module. (To verify that, for all \(f \in \mathfrak{r}_m\) , there exists \(c \neq 0\) in A such that \(cf \in \mathfrak{r}'\) , argue by induction on \(m\) , by writing

\[
f \left(\mathrm{X} _ {1}, \dots , \mathrm{X} _ {n}\right) = \mathrm{X} _ {1} f _ {1} \left(\mathrm{X} _ {1}, \dots , \mathrm{X} _ {n}\right) + \mathrm{X} _ {2} f _ {2} \left(\mathrm{X} _ {2}, \dots , \mathrm{X} _ {n}\right) + \dots + \mathrm{X} _ {n} f _ {n} \left(\mathrm{X} _ {n}\right)
\]

where the \(f_{j}\) are homogeneous of degree \(m - 1\) ; consider the polynomial \(g(\mathbf{X}_1, \ldots, \mathbf{X}_n) = \mathbf{X}_1 f_1(a_1, \ldots, a_n) + \mathbf{X}_2 f_2(a_2, \ldots, a_n) + \cdots + \mathbf{X}_n f_n(a_n)\) , observe that \(g \in \mathfrak{r}_1 \subset \mathfrak{r}'\) and form \(a_n^{m-1} f - \mathbf{X}_n^{m-1} g\) .

\begin{enumerate}
\def\labelenumi{(\alph{enumi})}
\setcounter{enumi}{2}
\tightlist
\item
  Deduce from (b) that when A is an integral domain, the kernel of the homomorphism r of (a) is the torsion module of S(a). Deduce that, for S(a) to be an integral domain, it is necessary and sufficient that S(a) be a torsion-free A-module. In particular, if the ideal a is a projective A-module, S(a) is an integral domain and isomorphic to R(a).
\end{enumerate}

\begin{enumerate}
\def\labelenumi{\arabic{enumi}.}
\setcounter{enumi}{4}
\tightlist
\item
  Let E be a vector space of finite dimension n over a field K.
\end{enumerate}

\begin{enumerate}
\def\labelenumi{(\alph{enumi})}
\tightlist
\item
  Show that for every integer \(m\) the symmetric power \(S^m(E)\) and the vector space \(S_m'(E)\) of symmetric contravariant tensors of order \(m\) have the same dimension \(\binom{m + n - 1}{n - 1} = \binom{n + m - 1}{m}\) .
\end{enumerate}

*(b) If K is of characteristic p \textgreater{} 0, show that the vector space \(\mathbb{S}_{p}^{\prime\prime}(\mathbf{E})\) of symmetrized tensors of order p has dimension \(\binom{n}{p}\) (observe that if, in a tensor product \(x_{1} \otimes x_{2} \otimes \cdots \otimes x_{p}\) of p vectors of E, \(x_{i} = x_{j}\) for an ordered pair of distinct indices i, j, then the symmetrization of this product is zero; use the fact that if \((\mathrm{H}_{i})_{1 \leqslant i \leqslant r}\) is a partition of \(\{1, 2, \ldots, p\}\) into non-empty sets, comprising at least two sets, the subgroup of \(S_{p}\) leaving each of the \(H_{i}\) stable has order not divisible by p). On the other hand, the canonical image in \(\mathbb{S}^{p}(\mathbf{E})\) of the space \(\mathbb{S}_{p}^{\prime}(\mathbf{E})\) of symmetric tensors of order p is of dimension n and is generated by the images of the tensors \(x \otimes x \otimes \cdots \otimes x\) (p times), where \(x \in E\) (use the same remark); the canonical mapping of \(\mathbb{S}_{p}^{\prime}(\mathbf{E})\) into \(\mathbb{S}^{p}(\mathbf{E})\) is therefore not bijective although the two vector spaces have the same dimension.*

